% GENERATED FILE. Edit canonical lessons, then run npm run generate:books.
% canonical-source-hash: fnv1a64:3903bc301d131798
% canonical-lessons: TA-C278-karumpu, TA-C278-ventaikkay, TA-C278-nilakkatalai, TA-C278-tarpucani, TA-C278-murukku

\chapter{Sugarcane, Okra, Peanut, Watermelon, Murukku}
\label{ch:ta-sugarcane-okra-peanut-watermelon-murukku}
\hlchaptermodality{278}

\begin{chapteropening}
\textbf{By the end of this chapter:} \emph{I can say sugarcane, okra, a peanut, a watermelon and murukku.}

Complete the last lesson of chapter 278: I can say sugarcane, okra, a peanut, a watermelon and murukku.
\end{chapteropening}

\section[\texorpdfstring{\ta{கரும்பு} (karumpu)}{கரும்பு (karumpu)}]{\ta{கரும்பு} (karumpu) — sugarcane}

\label{lesson:TA-C278-karumpu}



\begin{quote}
Before the new one: say the Tamil for a hobby, then the Tamil for an opinion.
\end{quote}

\subsection*{You'll want to know: \ta{கரும்பு}}
\textbf{\ta{கரும்பு}} — \emph{karumpu} — \textquotedblleft{}sugarcane\textquotedblright{}.

\begin{grammarlens}[title={What you've built}]
One more word for this chapter.
\end{grammarlens}

\subsection*{Your turn}
\begin{itemize}
\raggedright
  \item \emph{Say it:} \emph{karumpu}
  \item \emph{Say it:} \emph{karumpu}, once more
  \item \emph{Say it:} say \emph{karumpu}
  \item \emph{Recall:} say the Tamil for a hobby, then the Tamil for an opinion, then say \emph{karumpu} again
\end{itemize}

\begin{tcolorbox}[breakable,colback=teal!4,colframe=teal!35!black,title={Before you move on}]
What does \ta{கரும்பு} mean? (\textquotedblleft{}Sugarcane\textquotedblright{}.) Say it once more.
\end{tcolorbox}
\section[\texorpdfstring{\ta{வெண்டைக்காய்} (veṇṭaikkāy)}{வெண்டைக்காய் (veṇṭaikkāy)}]{\ta{வெண்டைக்காய்} (veṇṭaikkāy) — okra, ladies' finger}

\label{lesson:TA-C278-ventaikkay}



\begin{quote}
Before the new one: say the Tamil for an opinion, then the Tamil for sugarcane.
\end{quote}

\subsection*{You'll want to know: \ta{வெண்டைக்காய்}}
\textbf{\ta{வெண்டைக்காய்}} — \emph{veṇṭaikkāy} — \textquotedblleft{}okra, ladies' finger\textquotedblright{}.

\begin{grammarlens}[title={What you've built}]
One more word for this chapter.
\end{grammarlens}

\subsection*{Your turn}
\begin{itemize}
\raggedright
  \item \emph{Say it:} \emph{veṇṭaikkāy}
  \item \emph{Say it:} \emph{veṇṭaikkāy}, once more
  \item \emph{Say it:} say \emph{veṇṭaikkāy}
  \item \emph{Recall:} say the Tamil for an opinion, then the Tamil for sugarcane, then say \emph{veṇṭaikkāy} again
\end{itemize}

\begin{tcolorbox}[breakable,colback=teal!4,colframe=teal!35!black,title={Before you move on}]
What does \ta{வெண்டைக்காய்} mean? (\textquotedblleft{}Okra, ladies' finger\textquotedblright{}.) Say it once more.
\end{tcolorbox}
\section[\texorpdfstring{\ta{நிலக்கடலை} (nilakkaṭalai)}{நிலக்கடலை (nilakkaṭalai)}]{\ta{நிலக்கடலை} (nilakkaṭalai) — a peanut, a groundnut}

\label{lesson:TA-C278-nilakkatalai}



\begin{quote}
Before the new one: say the Tamil for sugarcane, then the Tamil for okra.
\end{quote}

\subsection*{You'll want to know: \ta{நிலக்கடலை}}
\textbf{\ta{நிலக்கடலை}} — \emph{nilakkaṭalai} — \textquotedblleft{}a peanut, a groundnut\textquotedblright{}.

\begin{grammarlens}[title={What you've built}]
One more word for this chapter.
\end{grammarlens}

\subsection*{Your turn}
\begin{itemize}
\raggedright
  \item \emph{Say it:} \emph{nilakkaṭalai}
  \item \emph{Say it:} \emph{nilakkaṭalai}, once more
  \item \emph{Say it:} say \emph{nilakkaṭalai}
  \item \emph{Recall:} say the Tamil for sugarcane, then the Tamil for okra, then say \emph{nilakkaṭalai} again
\end{itemize}

\begin{tcolorbox}[breakable,colback=teal!4,colframe=teal!35!black,title={Before you move on}]
What does \ta{நிலக்கடலை} mean? (\textquotedblleft{}A peanut, a groundnut\textquotedblright{}.) Say it once more.
\end{tcolorbox}
\section[\texorpdfstring{\ta{தர்பூசணி} (tarpūcaṇi)}{தர்பூசணி (tarpūcaṇi)}]{\ta{தர்பூசணி} (tarpūcaṇi) — a watermelon}

\label{lesson:TA-C278-tarpucani}



\begin{quote}
Before the new one: say the Tamil for okra, then the Tamil for a peanut.
\end{quote}

\subsection*{You'll want to know: \ta{தர்பூசணி}}
\textbf{\ta{தர்பூசணி}} — \emph{tarpūcaṇi} — \textquotedblleft{}a watermelon\textquotedblright{}.

\begin{grammarlens}[title={What you've built}]
One more word for this chapter.
\end{grammarlens}

\subsection*{Your turn}
\begin{itemize}
\raggedright
  \item \emph{Say it:} \emph{tarpūcaṇi}
  \item \emph{Say it:} \emph{tarpūcaṇi}, once more
  \item \emph{Say it:} say \emph{tarpūcaṇi}
  \item \emph{Recall:} say the Tamil for okra, then the Tamil for a peanut, then say \emph{tarpūcaṇi} again
\end{itemize}

\begin{tcolorbox}[breakable,colback=teal!4,colframe=teal!35!black,title={Before you move on}]
What does \ta{தர்பூசணி} mean? (\textquotedblleft{}A watermelon\textquotedblright{}.) Say it once more.
\end{tcolorbox}
\section[\texorpdfstring{\ta{முறுக்கு} (muṟukku)}{முறுக்கு (muṟukku)}]{\ta{முறுக்கு} (muṟukku) — murukku (a crunchy spiral snack)}

\label{lesson:TA-C278-murukku}



\begin{quote}
Before the new one: say the Tamil for a peanut, then the Tamil for a watermelon.
\end{quote}

\subsection*{You'll want to know: \ta{முறுக்கு}}
\textbf{\ta{முறுக்கு}} — \emph{muṟukku} — \textquotedblleft{}murukku (a crunchy spiral snack)\textquotedblright{}.

\begin{grammarlens}[title={What you've built}]
One more word for this chapter.
\end{grammarlens}

\subsection*{Your turn}
\begin{itemize}
\raggedright
  \item \emph{Say it:} \emph{muṟukku}
  \item \emph{Say it:} \emph{muṟukku}, once more
  \item \emph{Say it:} say \emph{muṟukku}
  \item \emph{Recall:} say the Tamil for a peanut, then the Tamil for a watermelon, then say \emph{muṟukku} again
\end{itemize}

\begin{tcolorbox}[breakable,colback=teal!4,colframe=teal!35!black,title={Before you move on}]
What does \ta{முறுக்கு} mean? (\textquotedblleft{}Murukku (a crunchy spiral snack)\textquotedblright{}.) Say it once more.
\end{tcolorbox}
